\chapter{Vorläufige Auswahl der Plattform}
\label{chap:auswahl}

Ziel des Vorprojekts ist keine endgültige Festlegung, sondern eine begründete Vorauswahl, die die spätere Umsetzung der Masterarbeit ermöglicht. Die Entscheidung stützt sich auf den Kriterienkatalog (Kapitel~\ref{sec:kriterienkatalog}) und die vergleichende Einordnung in Tabelle~\ref{tab:plattformvergleich-matrix}. 

\section{Gewichtung der Kriterien}
Für das Forschungsvorhaben stehen Reproduzierbarkeit und Wiederholbarkeit (K1), Automatisierung und Lifecycle-Steuerung (K2), Integrationsfähigkeit externer Werkzeuge (K3) sowie Beobachtbarkeit und Datenzugang (K4) im Vordergrund, da sie die Vergleichbarkeit von Experimenten und die spätere Evaluation von Cyber-Resilienz-Use-Cases direkt beeinflussen. Benutzer-, Rollen- und Rechteverwaltung (K6) sowie Wartbarkeit, Dokumentation und Community (K7) sind ebenfalls wesentlich, weil sie die langfristige Nutzbarkeit und den stabilen Betrieb einer Forschungsumgebung beeinflussen. Zugriffsmöglichkeiten (K5) werden unterstützend berücksichtigt, da sie die praktische Durchführbarkeit von Übungen und Experimenten maßgeblich erleichtern.

\section{Begründete Vorauswahl}
Aus Tabelle~\ref{tab:plattformvergleich-matrix} ergibt sich CyberRangeCZ als die im Vorprojekt am robustesten begründbare Option. Im Unterschied zu den Alternativen konnten zentrale Kriterien nicht nur anhand von Dokumentation, sondern zusätzlich durch den erfolgreich umgesetzten PoC praktisch untermauert werden. Damit liegen belastbare Hinweise vor, dass die Plattform den automatisierten Aufbau und Betrieb von Umgebungen (K2) sowie einen praktikablen Zugriff und die Durchführung von Trainingsinstanzen (K5) unterstützt. Darüber hinaus sprechen die aktuelle Weiterentwicklung, die verfügbare Dokumentation und Support-/Community-Strukturen (K7) sowie die vorhandenen Mechanismen für Mehrbenutzerbetrieb und Zugriffskontrolle (K6) für ein geringeres Implementierungs- und Betriebsrisiko im Verlauf der Masterarbeit.

KYPO wird als fachlich relevanter Vorgängeransatz eingeordnet, weist jedoch im Vergleich ein erhöhtes Nachhaltigkeitsrisiko auf, da es als Vorgängerprojekt gilt und der Schwerpunkt der Weiterentwicklung auf CyberRangeCZ liegt. OCR stellt zwar einen interessanten modularen Ansatz dar, konnte im Vorprojekt jedoch nicht in allen für die Auswahl entscheidenden Punkten praktisch verifiziert werden. Entsprechend wird OCR für die Masterarbeit vorerst nicht als primäre Zielplattform gewählt, kann aber als Referenz für sowie für eine mögliche spätere Erweiterung oder Gegenüberstellung herangezogen werden.

Auf Basis der Kriteriengewichtung und der im Vorprojekt gewonnenen Erkenntnis wird daher CyberRangeCZ als vorläufige Zielplattform für die Masterarbeit ausgewählt.
